\section{Containerisierung}

Um über Docker und Kubernetes sprechen zu können, müssen zunächst deren grundlegende Technik, die Container, besprochen werden. 
Es muss geklärt werden, was sie erfüllen, wie sie funktionieren und wie sie sich von virtuellen Maschinen (zur Vereinfachung ab sofort VM genannt) unterscheiden.

Container sollen die Aufgabe erfüllen, dass eine Anwendung auf jeglichen Geräten laufen kann, unabhängig von Betriebssystem, Umgebungsvariablen, Einstellungen und Ähnlichem. 
Dadurch soll erreicht werden, dass eine Anwendung nicht nur auf dem eigenen Gerät, sondern auf jedem Gerät funktioniert.

Dazu werden die Anwendungen in einem Container gestartet. Dieser führt die Anwendung in einer isolierten Umgebung aus. 

Dieses Konzept wurde, aber nicht von Containern eingeführt, sondern gab es schon mit VMs davor. 
Es unterscheidet sich aber wie VMs und Container das umsetzen, da eine VM die gesamt benötigte Hardware mit Softwarekomponenten simuliert,
während Container direkt die Infrastuktur des Host-System mit nutzt, wie den Kernel. \parencite[vgl. 54]{kofler2023docker}

Dadurch sind Container ressourcensparender als VMs. Zudem lassen sie sich schneller aufbauen.